\chapter*{ВВЕДЕНИЕ}                         % Заголовок
\addcontentsline{toc}{chapter}{ВВЕДЕНИЕ}    % Добавляем его в оглавление
Симуляция человеческого поведения является признанной практикой, 
как в социологии, так и в экономике{\cite{ziems2024can}. Однако, конвенциональные методы 
симуляции ограничены в моделировании всех существенных аспектов 
человеческих взаимодействий \cite{xie2024can}. Последние разработки в сфере искусственного 
интеллекта, в частности большие языковые модели (БЯМ, LLM), 
использовали инструменты агентного моделирования. \cite{openai2023gpt4,wu2023autogen, chen2023agentverse, park2023generative}. Большинство 
исследователей использовали LLM в различных игровых теоретических 
условиях, что является наиболее естественным способом исследования 
лежащих в основе экономических и социальных причин человеческого 
поведения. \cite{akata2023playing,horton2023large,aher2023using}.

Тем не менее, большинство существующих исследований исходят из 
недоказанной предпосылки, что агенты LLM ведут себя аналогично людям в 
моделировании. Более того, они игнорируют критический аспект 
эмоционального человеческого поведения, который существенно влияет на 
процесс принятия решений \cite{lerner2015emotion}. Чтобы устранить эти пробелы, была разработана 
новая методология для изучения соответствия между поведением LLM и 
человека на основе их сравнения в различных эмоциональных состояниях. В 
частности, мы сосредоточились на пяти базовых эмоциях: гнев, грусть, 
счастье, отвращение и страх, которые были выбраны на основе хорошо 
установленной классификации Пола Экмана ~\cite{ekman1999basic}. Мы предлагаем новый подход к 
наделению LLM эмоциями \cite{akata2023playing, brookins} с помощью эмоционального промптинга  \cite{li2023large}. 
исследуя его влияние на соответствие процессу принятия решений человека в 
играх с торгом.

Цель данной работы: исследование влияния эмпатии в больших 
языковых моделях при их использовании в процессе принятия решений в 
теории игр, разработка экспериментальной методологии для анализа влияния 
эмоций в играх с торгом, подбор проприетарных и свободных моделей 
искусственного интеллекта, проведение экспериментов в многоязычных 
средах с разными эмоциональными состояниями, анализ результатов.
